\documentclass[11pt,a4paper]{article}
\usepackage[T1]{fontenc}
\usepackage[utf8]{inputenc}
\usepackage[brazil]{babel}
\usepackage{lmodern}
\usepackage[margin=2cm]{geometry}
\usepackage{microtype,booktabs,longtable,array,amsmath,acro,indentfirst,float}
\setlength{\parindent}{0.5cm}
\usepackage[hidelinks]{hyperref}
\DeclareAcronym{cpu}{short=CPU,long=unidade central de processamento}
\DeclareAcronym{mlp}{short=MLP,long=perceptron multicamadas}
\DeclareAcronym{cnn}{short=CNN,long=rede neural convolucional}
\DeclareAcronym{sdr}{short=SDR,long=rádio definido por software}
\newcolumntype{L}[1]{>{\raggedright\arraybackslash}p{#1}}
\setlength{\tabcolsep}{3pt}
\renewcommand{\arraystretch}{1.18}
\setlength{\emergencystretch}{2em}
\linespread{0.98}
\title{Impactos de desempenho em inteligência artificial e rádio:\\ comparação entre ambientes com e sem Bitdefender}
\author{}
\date{}
\begin{document}\maketitle
As permissões disponíveis permitiram comparar dois computadores distintos. Os resultados são descritivos dos ambientes, com processadores, unidades e sistemas de arquivos diferentes.

O Bitdefender utilizou em média 1,96 processador lógico, com pico amostrado de 6,17, durante a nova campanha. A memória residente agregada média foi de 3,20 GiB. O componente ctc concentrou 72,55\% do tempo de CPU observado do agente.

Na rede convolucional PyTorch em CPU com dados gerados, a duração passou de 0,9666 s para 1,8065 s, aumento de 86,89\%. Na carga correspondente TensorFlow, houve redução de 47,51\%. Nas 1024 gravações e 1024 leituras de pequenos arquivos, a duração passou de 40,67 ms para 481,62 ms. Cada resultado reúne três repetições na mesma sessão.

Um atraso ilustrativo de 1 ms por lote, repetido em 437500 lotes e 100 épocas, soma 12,15 horas. No rádio a 8 milhões de amostras/s, 1 ms corresponde a 8000 amostras. O primeiro caso explica acumulação de custo e o segundo explica orçamento de tempo. Ambos são cenários hipotéticos para explicar as ordens de grandeza.

Os resultados comparam medições de 3 de novembro de 2026 no ambiente sem Bitdefender e de 5 de novembro no ambiente com Bitdefender.

\end{document}
